\section{An explicit seven-polynomial word and its pullbacks}
\label{app:seven-cubic-bank}

\begin{theorem}\label{thm:seven-cubic-bank}
For every integer $m\geq1$, there are a number field $K_m$, a set of
$14m$ distinct evaluation points, and a word whose complete list of
polynomials of degree at most $3m$ with at least $7m$ agreements consists
of exactly seven polynomials, each with exactly $7m$ agreements.
For each fixed $m$, the same holds over arbitrarily large prime fields.
For $m=1$ it holds over every prime field with $p\equiv1\pmod7$;
for $m=2$ it holds over every such field with $p>10^{10}$.
The latter case has length $28$, dimension $7$, and agreement $14$,
so has exactly quarter rate.
\end{theorem}

\begin{proof}
We first give a construction over any field $k$ of characteristic different
from $2,3,7$ containing a primitive seventh root $\zeta$.
Write
\[
 \eta=\zeta+\zeta^2+\zeta^4,\qquad
 \bar\eta=\zeta^3+\zeta^5+\zeta^6=-1-\eta.
\]
Then $\eta^2+\eta+2=0$ and $\eta\bar\eta=2$.  Set
\[
 A(X)=X^3-\eta X^2+\bar\eta X-1
     =\prod_{j\in\{1,2,4\}}(X-\zeta^j),\qquad
 P(X)=\eta X^3+X^2+X+\bar\eta,
\]
\[
 \alpha=\frac{\bar\eta}{\eta}=\frac{\eta-1}{2},
 \qquad c=\frac{3(\eta+3)}4.
\]
The construction rests on the two polynomial identities
\begin{align}
 P(X)-X^5&=(\bar\eta-X)(X-1)A(X),\label{eq:paley-seven-first}\\
 P(\alpha X)-cX^5&=-A(X)(cX^2+c\eta X+\bar\eta).
 \label{eq:paley-seven-second}
\end{align}
Both follow by multiplication and $\eta^2=-\eta-2$.
All parameters needed below are nonzero, and $\alpha^7\ne1$.
For completeness, under the involution $\eta\leftrightarrow\bar\eta$,
\[
 c\bar c=\frac92,\qquad
 (\alpha^7-1)(\bar\alpha^7-1)=\frac{7^3}{2^7}.
\]
These identities prove the assertions in every allowed characteristic,
including when the quadratic algebra splits.

Take the fourteen-point domain
$\Omega=\mu_7\sqcup\alpha\mu_7$ and define its received word by
\[
 w(\zeta^t)=\zeta^{5t},\qquad
 w(\alpha\zeta^t)=c\zeta^{5t}\quad(t\in\mathbb Z/7\mathbb Z).
\]
The seven candidates are
\[
 P_i(X)=\zeta^{5i}P(\zeta^{-i}X)
 \qquad(i\in\mathbb Z/7\mathbb Z).
\]
Their constant coefficients are distinct.  The identities above give four
prescribed matches for each candidate on $\mu_7$ and three on
$\alpha\mu_7$.  Each point of the first orbit has four prescribed matching
candidates, and each point of the second has three.  These incidences
already use the entire pairwise root budget:
\[
 7\binom42+7\binom32=63=3\binom72.
\]
Indeed, two distinct cubics share at most three evaluation values.
Consequently there are no further incidences: every candidate has exactly
seven matches, and the two orbit multiplicities are exactly four and three.

The same argument handles pullbacks and proves completeness.
Let $\psi(U)\in k[U]$ have degree $m$, and suppose each point of $\Omega$
has $m$ distinct preimages in $k$.  Pull back the domain, word, and seven
candidates along $\psi$.  Their degrees are at most $3m$, their agreement
counts are $7m$, and the matching multiplicities remain four and three.
If a different polynomial $Q$ of degree at most $3m$ has at least $7m$
agreements, counting its intersections with the seven candidates gives
\[
 3\,\#\{\text{agreements of }Q\}\leq7(3m).
\]
Equality is necessary, so every agreement lies above $\alpha\mu_7$.
Thus $Q$ agrees at all $7m$ such points with
$c\alpha^{-5}\psi(U)^5$.  Their difference has degree $5m$, with nonzero
leading coefficient, and cannot have $7m$ distinct roots.  This proves
completeness, even for candidates over an extension of $k$.

Over a number field containing $\zeta$, choose any shift $a$ outside
$\Omega$ and take $\psi(U)=U^m+a$.  Adjoining its finitely many fiber roots
gives the required $K_m$ and distinct fibers.  Reduction at sufficiently
large completely split primes of a normal closure preserves all these
identities and distinctness; there are infinitely many such primes.
For $m=1$, the construction is already defined over every
$\mathbb F_p$ with $p\equiv1\pmod7$.

For the asserted uniform prime bound when $m=2$, it remains to find a shift
$a\in\mathbb F_p$ such that all fourteen values $x-a$, $x\in\Omega$, are
nonzero squares.  Let $\chi$ be the quadratic character, with $\chi(0)=0$.
The distinct-root quadratic-character Weil bound
\cite[Lemma~2.1]{kim-yip-yoo-tuples} gives
\[
 \#\{a:\chi(x-a)=1\ \text{for every }x\in\Omega\}
 \geq \frac{p-98305\sqrt p}{16384}-14.
\]
To see the constant directly, expand
$2^{-14}\sum_a\prod_{x\in\Omega}(1+\chi(x-a))$.
A nonempty subset of size $j$ contributes error at most
$(j-1)\sqrt p$, and
$\sum_{j=1}^{14}\binom{14}{j}(j-1)=98305$.
Removing the fourteen possible zero locations costs at most fourteen.
The displayed lower bound is positive for $p>10^{10}$.
The map $\psi(U)=U^2+a$ therefore supplies the required distinct split
quadratic fibers.
\end{proof}

This is a fixed list-size construction.  It does not assert an unbounded
list, a superlinear number of exceptional challenges, or a universal
list bound of seven for arbitrary words at length twenty-eight.
